%
% frame_expiry.tex
%
% A Z specification of lux's frame time-to-live (TTL) expiry: a frame shown with
% a deadline is torn down once its deadline passes, unless a re-show refreshes it
% first. The one safety property that the TTL feature adds over the existing store
% is atomicity between a re-show and an expiry sweep: a frame re-armed with a fresh
% deadline must never be torn down by a stale (already-passed) deadline.
%
% Both the re-show (install the scene's roots AND arm the frame's deadline) and the
% sweep (decide a frame is due AND tear its roots down) run under the one store
% lock, so each is a single atomic step and they cannot interleave. This model
% encodes that atomicity and checks the safety invariant across every interleaving
% of re-show, clock advance, expiry, and manual close. The fidelity variant
% (Section~\ref{sec:fidelity}) splits the re-show into two steps -- dropping the
% atomicity the lock provides -- and ProB then reaches the exact bad state where a
% freshly re-shown frame has been torn down.
%
% Models the design in docs/architecture/target/target.md (the Hub is
% authoritative; a re-show is a whole-UI resend) as realised by
% src/punt_lux/domain/hub/frame_lifecycle.py (show/arm and expire/tear-down under
% one StoreLock) and src/punt_lux/domain/hub/frame_expiry.py (the deadlines).
%
% Type-check:  fuzz -t docs/frame_expiry.tex
% Model-check: make prob SPEC=docs/frame_expiry.tex   (coverage + deadlock)
%   Safety I1 is a reachability goal, not a state invariant (so the fidelity
%   variant can reach the bad state):
%     probcli docs/frame_expiry.tex -model_check \
%       -goal "#f.(f:FRAME & deadline(f)=dfuture & roots(f)=absent)" \
%       -p DEFAULT_SETSIZE 2
%   Atomic model: goal never satisfied. Fidelity variant (Section 8): satisfied.
%
% Formatting follows the fuzz house rules: schema lines under 80 columns,
% predicates broken at \land/\lor/\implies.
%

\documentclass[a4paper,10pt,fleqn]{article}
\usepackage[margin=1in]{geometry}
\usepackage{fuzz}
\usepackage[colorlinks=true,linkcolor=blue,citecolor=blue,urlcolor=blue]{hyperref}

\begin{document}

\title{lux Frame TTL Expiry: a Z Specification}
\author{Formal model of the re-show / expiry atomicity invariant}
\date{July 2026}
\maketitle

% ============================================================
\section{Introduction}
% ============================================================

A frame may carry a time-to-live. Shown with a positive TTL, it is removed once
that many seconds pass with no re-show refreshing it. The Hub is authoritative:
it arms the deadline when the frame is shown, sweeps the frames whose deadlines
have passed, and tears their scenes down through the same path a manual frame
close uses.

The sweep introduces a race the store did not have before. A frame's deadline can
pass at the very moment an agent re-shows that frame to refresh it. If the sweep
reads the passed deadline and tears the frame down \emph{after} the re-show has
re-installed its roots but \emph{before} the re-show has armed the fresh deadline,
the freshly shown frame is destroyed. The fix is that both the re-show and the
sweep run under the one store lock, so each is a single atomic step: the re-show
installs the roots and arms the new deadline together, and the sweep decides a
frame is due and tears it down together. Whichever runs second sees the other's
whole effect.

One safety property must hold across every interleaving:

\begin{description}
  \item[I1 --- a live deadline implies a shown frame.] A frame whose deadline is
    still in the future always has its roots installed. Equivalently: expiry never
    leaves a frame torn down while a fresh (future) deadline is armed on it. This is
    the property that fails if the re-show's ``install roots'' and ``arm deadline''
    are not one atomic step.
\end{description}

The model abstracts real time to three tokens --- a deadline is absent, in the
future, or already passed --- because I1 turns only on the \emph{order} in which a
frame's roots and its deadline change, not on any duration.

% ============================================================
\section{Basic Types}
% ============================================================

\begin{zed}
  [FRAME]
\end{zed}

\texttt{FRAME} is the set of frames the Hub can hold. One frame already exhibits
the race (the re-show and the sweep contend over a single frame's roots and
deadline); ProB's \texttt{DEFAULT\_SETSIZE} bounds the carrier and a second frame
confirms the invariant is per-frame.

Every operation acts on a frame input, so whenever any frame exists a
\texttt{ReShow} is enabled and the system can always act. The only state with
nothing enabled is the empty universe --- no frames at all, hence no work --- a
degenerate rest state, not a stuck system (see the deadlock note in the
verification section).

% ============================================================
\section{Free Types}
% ============================================================

\begin{zed}
  ROOTS ::= present | absent
\end{zed}

Whether a frame's scene roots are installed in the store. \texttt{present}: the
frame is shown. \texttt{absent}: torn down (never shown, closed, or expired).

\begin{zed}
  DEADLINE ::= dnone | dfuture | dpast
\end{zed}

A frame's TTL deadline, abstracted relative to ``now''. \texttt{dnone}: no TTL is
armed (a permanent frame). \texttt{dfuture}: a TTL is armed and has not yet passed.
\texttt{dpast}: the armed TTL has passed --- the frame is due for the sweep.

% ============================================================
\section{State}
% ============================================================

The state is flat: one schema. Its predicate carries only \emph{structural}
coherence --- \texttt{roots} and \texttt{deadline} are total over the frames. The
safety property I1 is deliberately \emph{absent} from the schema: were it a state
invariant, ProB would treat any operation that broke it as \emph{disabled} rather
than as a violation, and the fidelity check (Section~\ref{sec:fidelity}) could no
longer reach the bad state. I1 is therefore stated as a reachability goal and
checked against the reachable state space (see the verification section).

\begin{schema}{Frame}
  roots : FRAME \fun ROOTS \\
  deadline : FRAME \fun DEADLINE
\end{schema}

\texttt{roots} records, per frame, whether its scene is installed;
\texttt{deadline} records its TTL token. I1 --- any frame with a future deadline is
shown --- holds because the only operation that arms a future deadline
(\texttt{ReShow}) installs the roots in the same step, and the only operations that
remove the roots (\texttt{Expire}, \texttt{ManualClose}) clear the deadline in the
same step. The verification proves this is more than an assertion: the negating
state is unreachable.

% ============================================================
\section{Initialization}
% ============================================================

\begin{schema}{Init}
  Frame'
\where
  roots' = (\lambda f : FRAME @ absent) \\
  deadline' = (\lambda f : FRAME @ dnone)
\end{schema}

Every frame starts torn down with no deadline armed --- I1 holds vacuously.

% ============================================================
\section{Operations}
% ============================================================

\texttt{ReShow} is the whole show/re-show step: it installs the frame's roots and
arms a fresh (future) deadline \emph{together}, under the store lock. This is the
atomicity that guarantees I1: roots and a future deadline are set in one
indivisible transition, so no state ever has a future deadline without roots.

\begin{schema}{ReShow}
  \Delta Frame \\
  f? : FRAME
\where
  roots' = roots \oplus \{ f? \mapsto present \} \\
  deadline' = deadline \oplus \{ f? \mapsto dfuture \}
\end{schema}

\texttt{ReShowNoTtl} is a re-show that carries no TTL: it installs the roots and
clears any prior deadline, making the frame permanent. It cannot break I1 --- it
never arms a future deadline.

\begin{schema}{ReShowNoTtl}
  \Delta Frame \\
  f? : FRAME
\where
  roots' = roots \oplus \{ f? \mapsto present \} \\
  deadline' = deadline \oplus \{ f? \mapsto dnone \}
\end{schema}

\texttt{TimePass} is the environment advancing the clock past an armed deadline:
a frame's future deadline becomes passed. The roots are untouched --- the frame is
still shown, now merely due --- so I1 is preserved (the deadline is no longer
\texttt{dfuture}).

\begin{schema}{TimePass}
  \Delta Frame \\
  f? : FRAME
\where
  deadline~f? = dfuture \\
  deadline' = deadline \oplus \{ f? \mapsto dpast \} \\
  roots' = roots
\end{schema}

\texttt{Expire} is the sweep retiring a due frame: a frame whose deadline has
passed has its roots torn down and its deadline cleared, together, under the store
lock. Guarded on \texttt{dpast}, so the sweep only ever removes a frame whose TTL
has genuinely passed --- never one bearing a future deadline.

\begin{schema}{Expire}
  \Delta Frame \\
  f? : FRAME
\where
  deadline~f? = dpast \\
  roots' = roots \oplus \{ f? \mapsto absent \} \\
  deadline' = deadline \oplus \{ f? \mapsto dnone \}
\end{schema}

\texttt{ManualClose} is the user closing the frame (or a clear): its roots are torn
down and its deadline disarmed together, whatever the deadline was.

\begin{schema}{ManualClose}
  \Delta Frame \\
  f? : FRAME
\where
  roots' = roots \oplus \{ f? \mapsto absent \} \\
  deadline' = deadline \oplus \{ f? \mapsto dnone \}
\end{schema}

% ============================================================
\section{Verification}
% ============================================================

I1 is checked as a reachability goal: the negating state
\[
  B \defs \exists f : FRAME @ deadline~f = dfuture \land roots~f = absent
\]
--- a frame with a future deadline but no roots --- must be unreachable. In ProB:

\begin{verbatim}
probcli docs/frame_expiry.tex -model_check \
  -goal "#f.(f:FRAME & deadline(f)=dfuture & roots(f)=absent)" \
  -p DEFAULT_SETSIZE 2
\end{verbatim}

Against the atomic model this reports the goal is \emph{never} satisfied over the
whole reachable state space --- no reachable state has a frame with a future
deadline and absent roots, so no \texttt{Expire} ever fires on a frame a re-show
has just refreshed. The bare \texttt{-model\_check} additionally confirms every
operation is covered and the structural state predicate is preserved.

\paragraph{Deadlock.} There is nothing here to deadlock: the store is guarded by a
single reentrant lock, so there is no second lock and no acquisition order that
could cycle. Whenever a frame exists, \texttt{ReShow} is enabled, so the reachable
state space (which \texttt{-model\_check} traverses from \texttt{Init} over a
non-empty bounded carrier) has no stuck state. The constraint-based
\texttt{-cbc\_deadlock} does report one state with no enabled operation --- the
empty universe, $roots = \emptyset \land deadline = \emptyset$, where no frame
exists to act on. That state is unreachable from \texttt{Init} (which populates
every frame) and is not a concurrency deadlock; it is the trivial ``no frames, no
work'' rest state. It is recorded here rather than papered over.

% ============================================================
\section{Fidelity: dropping the atomicity reproduces the bug}
\label{sec:fidelity}
% ============================================================

The atomicity is load-bearing, and the model reproduces the defect when it is
removed. To confirm, split the single \texttt{ReShow} step into two
independently-schedulable steps --- as a design that armed the TTL in a
\emph{separate} store-lock acquisition from \texttt{show\_scene} would have:
\texttt{ReShowRootsBAD}, which sets $roots' = roots \oplus \{ f? \mapsto present \}$
and leaves \texttt{deadline} unchanged, and \texttt{ArmDeadlineBAD}, which sets
$deadline' = deadline \oplus \{ f? \mapsto dfuture \}$ and leaves \texttt{roots}
unchanged. These two schemas are deliberately \emph{absent} from the committed model above;
they are run only by temporarily swapping them in for \texttt{ReShow} and re-running
the goal check.

With them in place the goal $B$ becomes reachable, and ProB prints the trace. The
essence is that once arming a deadline is a step separate from installing the
roots, the two can be pulled apart: a frame can carry a future deadline while its
roots are absent. ProB's shortest witness arms a deadline on a frame that was
never shown (\texttt{ArmDeadlineBAD} straight from \texttt{Init}); the realistic
form is the re-show/expiry interleaving the lock is meant to forbid --- a re-show
installs the roots, the sweep expires the frame's already-passed deadline and tears
those roots down, and the re-show then arms its fresh deadline on the now
torn-down frame. Every such trace ends in $B$: a future deadline tracked for a
frame the Hub is not showing. The single atomic \texttt{ReShow} installs the roots
and arms the deadline in one indivisible step, so no trace can separate them, which
is why $B$ is unreachable in the committed model (all 16 reachable states visited,
goal never satisfied).

\end{document}
